

\begin{figure}
    \centering
    \includegraphics[width=\linewidth]{test_loss2.pdf}
    \caption{Top-1 test error plot for CIFAR100 (left) and ImageNet (right) classification. 
    Cutmix achieves lower test errors than the baseline at the end of training.
    }
    \label{fig:cls_loss}
\end{figure}


\subsection{Image Classification}
\label{sec:image_cls}

\subsubsection{ImageNet Classification}
\label{sec:imagenet_cls}
We evaluate on ImageNet-1K benchmark~\cite{ImageNet}, the dataset containing 1.2M training images and 50K validation images of 1K categories.
For fair comparison, we use the standard augmentation setting for ImageNet dataset such as re-sizing, cropping, and flipping, as done in~\cite{pyramidnet,ghiasi2018dropblock,densenet,GoogleNet}. 
We found that regularization methods {including} Stochastic Depth~\cite{stochasticdepth},  Cutout~\cite{devries2017cutout}, Mixup~\cite{zhang2017mixup}, and CutMix require a greater number of training epochs till convergence. 
Therefore, we have trained all the models for $300$ epochs with initial learning rate $0.1$ decayed by factor $0.1$ at epochs $75$, $150$, and $225$. 
The batch size is set to $256$. 
The hyper-parameter $\alpha$ is set to $1$. 
We report the best performances of CutMix and other baselines during training. 


\begin{table}[]
\small
\centering
\tabcolsep=0.1cm
\begin{tabular}{@{}lccc@{}}
\toprule
\multirow{2}{*}{Model}            & \multirow{2}{*}{\# Params} & \multirow{2}{*}{\begin{tabular}[c]{@{}c@{}}Top-1 \\ Err (\%)\end{tabular}} & \multirow{2}{*}{\begin{tabular}[c]{@{}c@{}}Top-5 \\ Err (\%)\end{tabular}} \\
                                  &                            &                        &        \\ \midrule
ResNet-152*                                     & 60.3 M     & 21.69    & 5.94\\
ResNet-101 + SE Layer*~\cite{SENet}             & 49.4 M     & 20.94    & 5.50\\
ResNet-101 + GE Layer*~\cite{GENet}             & 58.4 M     & 20.74    & 5.29\\
ResNet-50 + SE Layer*~\cite{SENet}              & 28.1 M     & 22.12    & 5.99\\
ResNet-50 + GE Layer*~\cite{GENet}              & 33.7 M     & 21.88    & 5.80\\ \midrule
ResNet-50 (Baseline)                            & 25.6 M     & 23.68    & 7.05\\
ResNet-50 + Cutout~\cite{devries2017cutout}     & 25.6 M     & 22.93    & 6.66\\
ResNet-50 + StochDepth~\cite{stochasticdepth}   & 25.6 M     & 22.46    & 6.27       \\
ResNet-50 + Mixup~\cite{zhang2017mixup}         & 25.6 M     & 22.58    & 6.40\\
ResNet-50 + Manifold Mixup~\cite{verma2018manifoldmixup}& 25.6 M     & 22.50     & 6.21\\
ResNet-50 + DropBlock*~\cite{ghiasi2018dropblock}   & 25.6 M    & 21.87     & {5.98}\\
ResNet-50 + Feature CutMix                      & 25.6 M    & 21.80  & 6.06\\ 
ResNet-50 + CutMix                              & 25.6 M    & \textbf{{21.40}}    & \textbf{{5.92}}\\  \midrule
\end{tabular}
\vspace{-0.1cm}
\caption{ImageNet classification results based on ResNet-50 model.
`*' denotes results reported in the original papers.
}
\vspace{-0.2cm}
\label{table:imagenet-res50}
\end{table}

We briefly describe the settings for baseline augmentation schemes.
We set the dropping rate of residual blocks to $0.25$ for the best performance of Stochastic Depth~\cite{stochasticdepth}.
The mask size for Cutout~\cite{devries2017cutout} is set to $112\times112$ and the location for dropping out is uniformly sampled. 
The performance of DropBlock~\cite{ghiasi2018dropblock} is from the original paper and the difference from our setting is the training epochs which is set to $270$.
Manifold Mixup~\cite{verma2018manifoldmixup} applies Mixup operation on the randomly chosen internal feature map.
We have tried $\alpha=0.5$ and $1.0$ for Mixup and Manifold Mixup and have chosen $1.0$ which has shown better performances.
It is also possible to extend CutMix to feature-level augmentation (Feature CutMix).
Feature CutMix applies CutMix at a randomly chosen layer per minibatch as Manifold Mixup does.


\begin{table}[]
\small
\centering
\begin{tabular}{@{}lccc@{}}
\toprule
\multirow{2}{*}{Model}     & \multirow{2}{*}{\# Params} & \multirow{2}{*}{\begin{tabular}[c]{@{}c@{}}Top-1 \\ Err (\%)\end{tabular}} & \multirow{2}{*}{\begin{tabular}[c]{@{}c@{}}Top-5 \\ Err (\%)\end{tabular}} \\
                          &                            &                               &  \\ \midrule
ResNet-101 (Baseline)~\cite{resnet}    & 44.6 M    & 21.87     & 6.29\\
ResNet-101 + Cutout~\cite{devries2017cutout} & 44.6 M    & 20.72    & 5.51\\ 
ResNet-101 + Mixup~\cite{zhang2017mixup}    & 44.6 M    & 20.52    & 5.28\\ 
ResNet-101 + CutMix                     & 44.6 M    & \textbf{20.17}    & \textbf{5.24}\\ \midrule
ResNeXt-101 (Baseline)~\cite{resnext}   & 44.1 M    &    21.18   &  5.57 \\
ResNeXt-101 + CutMix                    &  44.1 M    &  \textbf{19.47}  & \textbf{5.03}  \\ 
\toprule
\end{tabular}
\vspace{-0.2cm}
\caption{Impact of CutMix on ImageNet classification for ResNet-101 and ResNext-101.}
\vspace{-0.4cm}
\label{table:imagenet-deeper}
\end{table}



\noindent\textbf{Comparison against baseline augmentations: }
Results are given in Table~\ref{table:imagenet-res50}.
We observe that CutMix achieves the best result, {$\textbf{21.40\%}$} top-1 error, among the considered augmentation strategies.
CutMix outperforms Cutout and Mixup, the two closest approaches to ours, by $+1.53\%$ and $+1.18\%$, respectively.
On the feature level as well, we find CutMix preferable to Mixup, with top-1 errors $21.78\%$ and $22.50\%$, respectively.

\noindent\textbf{Comparison against architectural improvements: }
We have also compared improvements due to CutMix versus architectural improvements (\eg greater depth or additional modules). 
We observe that CutMix improves the performance by $\textbf{+2.28\%}$ while increased depth (ResNet-50 $\rightarrow$ ResNet-152) boosts $+1.99\%$ and SE~\cite{SENet} and GE~\cite{GENet} boosts $+1.56\%$ and $+1.80\%$, respectively.
Note that unlike above architectural boosts improvements due to CutMix come at little or memory or computational time.

\noindent\textbf{CutMix for Deeper Models: }
We have explored the performance of CutMix for the deeper networks, ResNet-101~\cite{resnet} and ResNeXt-101 (32$\times$4d)~\cite{resnext}, on ImageNet.
As seen in Table~\ref{table:imagenet-deeper}, we observe $ \textbf{+1.60\%}$ and $\textbf{+1.71\%}$ respective improvements in top-1 errors due to CutMix.
